% Direct-authority diplomatic French for printed p.190.
\clearpage
\SourcePage{190}

\LayerTarget{tohoku:corollary:4.2.4:1}
\LayerTarget{tohoku:eq:4.2.10}
\[
\tag{5.2.10}
\mathbf{Ext}^n(A,B)=H^n(\mathbf{Hom}_{\mathbf O}(\mathbf{L},B))\qquad\text{\textit{d'où}}
\]
\LayerTarget{tohoku:eq:4.2.10-bis}
\[
\tag{5.2.10 \textit{bis}}
\mathbf{Ext}^n(A,B)(x)=
\mathrm{Ext}_{{\mathbf O}(x)}^n(A(x),B(x)),
\]

\LayerTarget{tohoku:proof:corollary-4.2.4:1}
Il suffit en effet d'appliquer prop. 2.5.1.

Appliquons maintenant prop. 2.5.4, on trouve :

\LayerTarget{tohoku:corollary:4.2.4:2}
\textsc{Corollaire 2.} \textit{Sous les conditions précédentes, le
foncteur spectral du th. \textup{4.2.1} est isomorphe au foncteur spectral
$\mathrm{II}\Gamma(\mathbf{Hom}_{\mathbf O}(\mathbf{L},B))$, en particulier on a}
\[
\mathrm{Ext}_{\mathbf O}^n(X;A,B)=
\mathfrak{R}^n\Gamma\mathbf{Hom}_{\mathbf O}(\mathbf{L},B))
\]
\textit{\textup{(}où le deuxième membre est le n.ème groupe de hyperhomologie de
$\Gamma$ par rapport au complexe $\mathbf{Hom}_{\mathbf O}(\mathbf{L},B)$, cf.\ \textup{2.4.).}}

\LayerTarget{tohoku:construction:canonical-resolvent-spectral-sequence}
Comme nous disposons aussi d'un foncteur résolvant pour $\Gamma$, par exemple
le foncteur $\Gamma {\mathrm C}(F)$ (où ${\mathrm C}(F)$ est la \textit{résolution canonique} de
$F$, cf. 2.3), la suite spectrale
$\mathrm{II}\Gamma(\mathbf{Hom}_{\mathbf O}(\mathbf{L},B))$ s'explicite, c'est la première suite
spectrale du bicomplexe $\Gamma {\mathrm C}(\mathbf{Hom}_{\mathbf O}(\mathbf{L},B))$ (où on a pris
comme premier degré celui qui provient de ${\mathrm C}$). Or on a un isomorphisme
fonctoriel
\LayerTarget{tohoku:eq:4.2.11}
\[
\tag{4.2.11}
{\mathrm C}(\mathbf{Hom}_{\mathbf O}(\mathbf{L},B))=
\mathbf{Hom}_{\mathbf O}(\mathbf{L},{\mathrm C}(B)).
\]
En effet, on a pour tout ${\mathbf O}$-Module $L$ un homomorphisme (fonctoriel)
\LayerTarget{tohoku:eq:4.2.12}
\[
\tag{4.2.12}
{\mathrm C}(\mathbf{Hom}_{\mathbf O}(L,B))\to\mathbf{Hom}_{\mathbf O}(L,{\mathrm C}(B))
\]
\LayerTarget{tohoku:proof:equations-4.2.11-4.2.12}
(où au deuxième membre, ${\mathrm C}(B)$ est considéré de façon évidente comme un
${\mathbf O}$-Module), défini par récurrence sur la dimension des composantes de ${\mathrm C}$ à
partir de l'homomorphisme naturel
${\mathrm C}^0(\mathbf{Hom}_{\mathbf O}(L,B))\to\mathbf{Hom}_{\mathbf O}(L,{\mathrm C}^0(B))$, déduit des
homomorphismes naturels
$\mathbf{Hom}_{\mathbf O}(L,B)(x)\to
\mathrm{Hom}_{{\mathbf O}(x)}(L(x),B(x))$. Ce dernier est un isomorphisme si $L$ est
pseudo-cohérent (prop. 4.1.1), d'où on conclut que dans ce cas, (4.2.12) est un
isomorphisme, ce qui établit en particulier la formule 4.2.11. On en conclut
$\Gamma {\mathrm C}(\mathbf{Hom}_{\mathbf O}(\mathbf{L},B))=
\mathrm{Hom}_{\mathbf O}(\mathbf{L},{\mathrm C}(B))$, d'où le

\LayerTarget{tohoku:corollary:4.2.4:3}
\textsc{Corollaire 3.} \textit{$\mathbf{L}$ étant comme dans la
proposition \textup{4.2.1}, le foncteur spectral du théorème \textup{4.2.1} est donné par la
première suite spectrale
$I\mathrm{Hom}_{\mathbf O}(\mathbf{L},{\mathrm C}(B))$ du bicomplexe
$\mathrm{Hom}_{\mathbf O}(\mathbf{L},{\mathrm C}(B))$, où ${\mathrm C}(B)$ est la résolution canonique
\textup{(}cf.\ \textup{2.3)} du faisceau abélien $B$ \textup{(}dont la graduation est prise comme premier
degré du bicomplexe\textup{)}. En particulier, on a}
\[
\mathrm{Ext}_{\mathbf O}^n(X;A,B)=
H^n(\mathrm{Hom}_{\mathbf O}(\mathbf{L},{\mathrm C}(B)).
\]

\LayerTarget{tohoku:remark:ext-with-supports}
\textsc{Remarques.} 1) On peut aussi définir les foncteurs
$\mathrm{Ext}_{{\mathbf O},\varphi}^p(X;A,B)$ comme les foncteurs dérivés droits, par
rapport à $B$, du foncteur
$\mathrm{Hom}_{{\mathbf O},\varphi}(A,B)=
\Gamma(\mathbf{Hom}_{\mathbf O}(A,B))$ ($\varphi$ étant un antifiltre de parties
fermées de $X$). Les développements précédents sont encore valables. Nous
n'aurons pas à nous en servir.

\LayerTarget{tohoku:remark:finiteness-in-theorem-4.2.2}
2) Il est essentiel dans le théorème 4.2.2 que $A$ satisfasse à une condition
de finitude. Si par exemple $A={\mathbf O}^{(I)}$ oû $I$ est infini, alors le deuxième
membre de (4.2.8) est nul pour tout $x$ et tout $p>0$, cependant on n'aura
\textit{pas} en général $\mathbf{Ext}_{\mathbf O}^1(A,B)=0$, car le foncteur
$B\to\mathbf{Hom}_{\mathbf O}({\mathbf O}^{(I)},B)=B^I$ n'est en général pas un foncteur
exact.

\TohokuContinuation{191}
